\chapter{Conclusion and Outlook}
\label{chap:conclusion}

Functional linear regression turns a familiar linear prediction problem into
an inverse problem.  The covariance operator suppresses weak directions, its
empirical version is finite rank, and any attempt to invert it must decide both
which information to retain and how to represent the resulting computation.
This thesis has examined those decisions through spectral truncation, raw and
orthogonalised Krylov least squares, and direct conjugate-gradient iteration.
Its main conclusion is that regularisation cannot be separated from purpose or
from numerical realisation.  A rule that is effective for slope estimation can
be inadequate for prediction or inference, and two procedures that target the
same exact-arithmetic subspace can have radically different reliability in
floating-point coordinates.

The argument combines operator geometry, fully specified algorithms, controlled
Monte Carlo experiments, a separate inference study, and one observed
spectroscopic application.  The evidence is deliberately conditional.  It
does not produce a universally best method; it identifies when target
geometry, stopping, and numerical representation change the conclusion.

\section{Answers to the research questions}
\label{sec:conclusion-research-questions}

\subsection{Spectral and iterative regularisation}

The first question asked how spectral and iterative regularisation compare as
the covariance structure and slope change.  The simulations show that the
answer depends on the loss used to define success.  CG--FPLS stops very early
under the implemented discrepancy rule.  That aggressive regularisation often
reduces variance in the ordinary $L^2$ slope norm, giving it the smallest mean
ISE in Models~2 and 3 and a concentrated slope-error distribution.  The same
stopping can leave response-relevant directions unresolved.  Its mean
prediction error exceeds that of Arnoldi FPLS in all three models, most clearly
in Model~2, where their mean MSPEs are 1.999 and 1.442.

FPCR and the response-informed Krylov procedures order information differently.
FPCR follows predictor variance, whereas FPLS builds spaces from both
$\widehat K$ and $\widehat r$.  The response-informed methods can therefore
reach a useful predictive relation in relatively few directions, but no fixed
component count has the same meaning across the two geometries.  In Models~1
and 3, FPCR and Arnoldi have almost identical mean prediction errors; in
Model~2, Arnoldi has the smallest stable-method mean and median MSPE.  The
simulation result is thus a target-dependent trade-off rather than a categorical
advantage of spectral or iterative regularisation.

The tablet application repeats the prediction-side cost of the implemented CG
rule, which stops at three steps: CG--FPLS has benchmark RMSE 5.293, compared
with 4.607 for Raw and 4.643 for Arnoldi.  Because the physical slope is
unknown, the application cannot determine whether early stopping improves
ordinary slope recovery.  It confirms only that, for this assay prediction
problem, the remaining response information matters enough to improve
benchmark prediction.

\subsection{Finite-precision Krylov representation}

The second question concerned the reliability of theoretically equivalent
FPLS formulations.  At a common valid dimension, raw and Arnoldi FPLS span the
same Krylov space and target the same restricted response least-squares fit.
Their paired stable-subset simulation behaviour is correspondingly almost
identical.  The equivalence fails as an implementation guarantee, however,
because the raw powers become nearly dependent.  Raw FPLS is flagged in
0.40\%, 8.18\%, and 0.32\% of the 5,000 simulations in Models~1--3.  Those
rare fits create catastrophic ISE and MSPE tails that are absent after
orthogonalisation.  Medians alone conceal this reliability difference.

The observed spectra make the distinction even sharper.  Twenty-four of 25
raw development fits and the final raw fit are flagged.  Its selected
dimensions range from four to 70, and the final $m=55$ basis and design
condition numbers reach $4.50\times10^{15}$.  Arnoldi selections stay between
four and six; the final $m=5$ basis condition is one and its orthogonality
defect is $7.4\times10^{-16}$.  Yet the two final prediction vectors correlate
at 0.999934 and differ by only 0.174 assay units in RMS.  Their benchmark RMSE
difference is just 0.036.

This is the clearest answer to the second question.  Orthogonalisation does not
promise a lower loss in each realised dataset, nor does it define a different
exact fixed-$m$ target.  It keeps the intended Krylov fit accessible and its
selection path interpretable when the raw power coordinates cease to be
trustworthy.  A finite prediction score cannot by itself certify numerical
reliability.

\subsection{Stopping for estimation and inference}

The third question asked why an inference-compatible CG index can differ from
one selected for estimation.  The distinction follows from the quantity that
regularisation is allowed to leave behind.  Estimation uses early stopping to
control amplification of sampling noise.  The direct-moment test instead
requires the residual
$\widehat r-\widehat K\widehat\beta_m$ to be negligible on the root-$n$
sampling scale.  Crossing a fixed estimation-oriented discrepancy threshold
does not imply that stronger requirement.

The stopping experiment illustrates the consequence with an
estimation-oriented variant.  It uses the earlier discrepancy formula and
constants, but caps the variance pilot at 20 components rather than 70 and
uses the stabilised moment-residual path rather than the short recurrence.
It therefore does not reproduce the complete estimation procedure unchanged.
This comparator selects one or two components in most samples, satisfies its
own realised threshold, yet produces 5\%-level null rejection frequencies of 52.70\%,
96.02\%, and 71.94\% in Models~1--3.  Its finite-iteration residual remains
large on the root-$n$ scale.  Continuing the path reduces both the
scaled residual and the difference from the direct-moment statistic.  At
$m=70$, the two versions make identical stored decisions, with rejection
frequencies of 6.20\%, 6.00\%, and 5.26\% in that stopping-diagnostic
experiment.  In the separate principal size stream, the direct-moment
oracle-calibrated rejection frequencies are 5.90\%, 5.84\%, and 6.10\%.  Both
sets show modest liberal behaviour at $n=100$.  Once the iterative residual is
negligible, the discrepancy from the continuous-design oracle reference
combines finite-sample calibration effects, the fixed-grid approximation, and
simulation error in the reference distribution and critical value.

The inference study also establishes the narrower positive results that the
design supports.  Empirically calibrated direct-moment power rises sharply
when $n$ increases from 100 to 200; the finite-dimensional inversions are
non-empty, bounded, and full rank in every replication.  Neither result makes
$m=70$ a universal stopping rule.  The general lesson is to diagnose the
root-$n$ moment residual and to validate a feasible calibration, rather than
reuse an estimation index for formal inference.

\subsection{Regularisation in observed spectroscopy}

The fourth question asked whether the same trade-offs appear with observed
functional predictors.  All four procedures recover substantial assay signal
on the supplied benchmark, with $R^2$ from 0.887 to 0.914.  Raw and Arnoldi
form the leading predictive pair.  The paired iid interval for Raw-minus-
Arnoldi RMSE is $[-0.057,-0.017]$, but the point difference is less than 1\%
of Arnoldi's RMSE and their tablet-level errors move almost together.  It is a
conditional difference that is small relative to the benchmark prediction
errors.  Its practical importance cannot be determined without an assay scale
and an application-specific tolerance.

The stable empirical conclusion is therefore not that Raw wins.  Arnoldi
realises essentially the same predictive relation with a five-component,
orthonormal representation, while the nominal 55-component raw fit is severely
ill-conditioned.  The primary disadvantage of CG is consistent with the cost
of aggressive stopping seen in simulation, and FPCR's larger selected index
illustrates the different ordering imposed by spectral truncation.  Trimming,
smoothing, and changing instrument alter both the errors and selected
dimensions; the Raw--Arnoldi point ordering itself changes across variants.
The application transfers the mechanisms developed in the controlled study
without converting them into a universal ranking.

\section{Integrated implications for practice}
\label{sec:conclusion-implications}

The first practical implication is to specify the target before choosing a
regularisation rule.  Ordinary slope error, covariance-weighted prediction,
and a moment-based inferential approximation reward different compromises.
An early index that suppresses slope variance can still omit predictive signal,
while a useful prediction index can leave an inferential residual that is much
too large.  Reporting a single generic measure of fit does not resolve this
distinction.

Second, the basis is part of the procedure in finite precision.  Krylov-space
notation identifies a mathematical subspace, but software must still produce
coordinates in that space and solve the induced least-squares problem.  Raw
powers can be adequate at low dimension and fail abruptly later.  Arnoldi with
reorthogonalisation and QR supplies a defensible default when the aim is to
follow the APLS path, especially when its prediction is effectively tied with
the raw fit.  Condition numbers, orthogonality defects, effective dimensions,
and invalid candidates should accompany loss summaries.

Third, component counts must be interpreted within their algorithms.  A
principal component, a raw or Arnoldi Krylov direction, and a CG stopping
index do not constitute equal units of effective complexity.  The comparison
in this thesis intentionally preserves each procedure's native objective and
tuning rule.  It therefore evaluates complete procedures, not isolated
regularisers at an artificial common index.

Finally, reproducibility must extend below an aggregated results table.  The
estimation archive retains replication-level losses, selected dimensions,
instability diagnostics, and cached quantities supporting the distributional
figures.  The inference archive retains statistics, confidence-set summaries,
diagnostics, and configuration information.  The spectroscopy archive
additionally preserves observation-level predictions, fitted coefficients,
tuning paths, exact resampling indices, source snapshots, environment records,
and checksums.  These records support checks of the reported statistical and
numerical behaviour, with the scope of exact replay determined by what each
archive retains.

\section{Limitations}
\label{sec:conclusion-limitations}

The estimation Monte Carlo evidence covers three dense, error-free functional designs
with Gaussian scores, homoskedastic errors, $n=100$, and fixed implementation
rules.  These choices isolate mechanisms but do not span irregular sampling,
measurement error, heavy-tailed or dependent errors, sparse curves, or a broad
range of eigenvalue and slope classes.  Monte Carlo rankings are conditional
on the stated tuning procedures and should not be interpreted as minimax or
uniform comparisons.

The inference analysis is still further from a ready observational workflow.
Its principal calibration uses continuous-design population quantities known
from the simulation; their covariance weights differ slightly from those of
the fixed grid.  Its power calculations use simulation-specific empirical null
critical values.  Only one perturbation direction is examined, and the
confidence-set display restricts candidates to five basis coefficients.  No
stored result validates a feasible plug-in, multiplier, or bootstrap
calibration on observed data.  This is why the tablet analysis makes no formal
functional inference claim.

The empirical evidence comes from one response in one public archive.  There
is no independent tablet batch, laboratory, site, or acquisition-period test;
instrument~2 repeats measurements on the same tablets.  Benchmark responses
had already influenced dataset selection, so the supplied test split is not a
prospectively blind holdout.  Its bootstrap intervals hold all fitted models
fixed and omit variation from development sampling, preprocessing, and tuning.
No batch identifiers explain possible archive-order dependence, assay units
are not explicit in the MAT object, and the true physical slope is unavailable.
The coefficients consequently support replay and numerical comparison, not
causal chemical interpretation.  Finally, the publicly downloadable archive
contains no explicit open-data licence, limiting redistribution of the raw
file.

\section{Directions for further work}
\label{sec:conclusion-future-work}

The most immediate methodological task is feasible calibration of the late-
stopped statistic.  Plug-in, multiplier, or bootstrap approximations should be
studied under the same finite-sample discipline used here, with critical-value
uncertainty separated from residual iteration error.  A related goal is a
data-dependent inference stopping rule that directly controls the scaled
moment residual, rather than selecting a prediction index and hoping it is
late enough.

On the numerical side, stability-aware selection could combine predictive
loss with an explicit admissibility rule for conditioning and rank.  When
several candidates have practically indistinguishable validation error, such a
rule could prefer the simplest stable representation.  Comparisons with ridge,
spline, and other functional regularisers would show whether the target and
conditioning lessons persist beyond truncation and Krylov methods.

For spectroscopy, the decisive next step is genuine external validation across
tablet batches, instruments, laboratories, and acquisition periods.  Known
groups would permit batch-aware or order-aware resampling, while a bootstrap
that refits preprocessing, tuning, and models would address a broader
uncertainty target than the fixed-model intervals used here.  Measurement-
error models and irregular-grid methods would further connect the idealised
functional model to practical acquisition.  Such studies should freeze the
primary response and preprocessing rules before opening the external outcomes.

\section{Final perspective}
\label{sec:conclusion-final-perspective}

Regularisation in functional linear regression is simultaneously statistical,
computational, and purpose-dependent.  Early stopping can improve estimation
yet invalidate inference; theoretically equivalent Krylov spaces can behave
very differently in finite-precision coordinates; and a strong realised
prediction score need not certify a stable estimator.  Reliable practice
therefore requires matching the stopping rule to the target, using numerically
defensible representations, and preserving the diagnostics needed to
distinguish predictive success from computational fragility.
